% Publication-ready replacement for the unsupported CMP(1) paragraph in
% G18_FIXED_SPACING_CARRIER_BRIDGE_INSERT.tex, lines 103--184.
% This file is input by G18_FIXED_SPACING_CARRIER_BRIDGE_INSERT.tex.

\subsubsection{Coefficient close projection and transport to the physical GNS sector}
\label{sec:g18-cmp1-complete}

We use the three-outgoing-link cell decomposition and the gap-one
normalization of Proposition~\ref{prop:periodic-yarotsky-descent}.  If
$J\Subset\mathbb Z^3$ is a nonempty exact cell support, write
$\mathcal H'_J$ for the corresponding excited tensor factor and
$H_{J,0}$ for the restriction of the free cell Hamiltonian.  Thus
\begin{equation}
 H_{J,0}\ge |J|\,1_{\mathcal H'_J}.
 \label{eq:g18-cmp1-cell-gap}
\end{equation}
The superscript minus below means both local gauge invariance and odd charge
conjugation parity.  Define
\begin{align}
 \mathcal X_0^-&=
 \left\{v=(v_J):v_J\in\mathcal H'_J,
       \ \|v\|_0:=\sum_J\|v_J\|<\infty\right\}^-,
 \label{eq:g18-cmp1-x0}\\
 \mathcal X_1^-&=
 \left\{v\in\mathcal X_0^-:v_J\in\operatorname{Dom}H_{J,0},
       \ \|v\|_1:=\sum_J\|H_{J,0}v_J\|<\infty\right\}.
 \label{eq:g18-cmp1-x1}
\end{align}
The empty coefficient is absent in the odd sector.  Consequently
\eqref{eq:g18-cmp1-cell-gap} makes $\|\cdot\|_1$ a norm, and the finite
coefficient families with entries in $\operatorname{Dom}H_{J,0}$ form a
core, denoted $\mathcal c_{00}^-$, for
\begin{equation}
 D(v_J)_J=(H_{J,0}v_J)_J,
 \qquad \operatorname{Dom}D=\mathcal X_1^-.
 \label{eq:g18-cmp1-D}
\end{equation}

\begin{lemma}[Exact-support symmetry]
\label{lem:g18-cmp1-exact-support-symmetry}
The exact-support projections commute with every local gauge
transformation and with charge conjugation.  Hence $D$, Yarotsky's
coefficient perturbation $F_u$, the free shell projection $R^-$, and the
coefficient Riesz projection $\mathcal P_u^-$ preserve
$\mathcal X_j^-$, $j=0,1$.
\end{lemma}

\begin{proof}
Let $p_x=|\Omega_x\rangle\langle\Omega_x|$ and $q_x=1-p_x$ on a grouped
cell.  In a finite box the exact-support projector is a tensor product of
$q_x$ for $x\in J$ and $p_x$ for $x\notin J$.  A gauge transformation at a
vertex acts on incident link factors by left or right regular
representations.  After the links are regrouped into cells this action still
factorizes over the cell tensor factors, and every factor fixes the constant
cell vacuum $\Omega_x$.  It therefore commutes with $p_x$, $q_x$, and every
exact-support projector.  Charge conjugation is the linkwise pullback by the
group automorphism $U\mapsto\overline U$; it also factorizes and fixes
$\Omega_x$, so the same conclusion holds.

 The free electric Hamiltonian and the plaquette interaction commute with
 these symmetries.  The symmetries fix the product vacuum and preserve the
 normalization
 $\langle\widetilde\Omega_\Lambda(u),\Omega_{\Lambda,0}\rangle=1$.
 Since the interacting ground line is unique in the small-coupling domain,
 this normalization fixes its representative and makes that representative
 invariant.  Uniqueness of the finite-volume coefficient expansion then
 implies that $D$ and $F_u$
commute with the induced coefficient actions.  Passing to Yarotsky's
thermodynamic coefficient limits preserves these commutators.  Functional
calculus for $D$ and the contour definition of $\mathcal P_u^-$ prove the
last assertion.
\end{proof}

\begin{lemma}[Unweighted coefficient close projection]
\label{lem:g18-cmp1-close-projection-complete}
Let $g=\sup_x\|\widehat\phi_x(u)\|=27|u|$.  There is a function
$a(g)\to0$ such that
\begin{equation}
 \|F_uv\|_0\le a(g)\|v\|_1,
 \qquad v\in\mathcal X_1^-.
 \label{eq:g18-cmp1-relative-bound}
\end{equation}
 Assume in addition the spectral-enclosure conditions
 \begin{equation}
  27|u|<c_1,\qquad
  \lambda:=27c_2|u|<\frac1{33}.
  \label{eq:g18-cmp1-enclosure-smallness}
 \end{equation}
 On the circle
 \begin{equation}
  \gamma_u=\{z:|z-4|=r_u\},
  \qquad r_u=\frac{1-\lambda}{8},
  \label{eq:g18-cmp1-dynamic-contour}
 \end{equation}
 put
 \begin{equation}
  R^-=-\frac1{2\pi i}\oint_{\gamma_u}(D-z)^{-1}\,dz,
  \qquad
  \mathcal P_u^-=-\frac1{2\pi i}
        \oint_{\gamma_u}(D+F_u-z)^{-1}\,dz .
  \label{eq:g18-cmp1-projections}
 \end{equation}
 If $q(g):=35a(g)<1$, then
\begin{equation}
 \boxed{\quad
 \|\mathcal P_u^--R^-\|_{\mathcal X_0^-\to\mathcal X_0^-}
 \le \kappa(g):=\frac{q(g)}{1-q(g)}.
 \quad}
 \label{eq:g18-cmp1-kappa}
\end{equation}
In particular, $\kappa(g)\to0$; after decreasing the small-coupling
interval, $\kappa(g)<1$.
\end{lemma}

\begin{proof}
Fix any $0<\rho<1$.  Yarotsky's estimate~(14), for an input of exact
support $J$, gives
\[
 \sum_K\|(F_uv_J)_K\|\rho^{-(d(K,J)+1)}
 \le C_{14}(\rho)g|J|\|v_J\|.
\]
Since the weight on the left is at least $\rho^{-1}$,
\[
 \sum_K\|(F_uv_J)_K\|
 \le \rho C_{14}(\rho)g|J|\|v_J\|
 \le \rho C_{14}(\rho)g\|H_{J,0}v_J\|.
\]
Summing the columns proves~\eqref{eq:g18-cmp1-relative-bound} with
$a(g)=\rho C_{14}(\rho)g$.  This also proves that $D+F_u$, with domain
$\mathcal X_1^-$, is closed when $a(g)<1$.

 The retained-shell gaps put $\gamma_u$ in the resolvent set of $D$ and give
 \[
  K_0:=\sup_{z\in\gamma_u}\|(D-z)^{-1}\|_{0\to0}
  \le\frac8{1-\lambda}.
 \]
 Moreover,
 \[
  D(D-z)^{-1}=1+z(D-z)^{-1},
  \qquad
  K_1:=\sup_{z\in\gamma_u}\|D(D-z)^{-1}\|_{0\to0}
  \le2+\frac{32}{1-\lambda}<35,
 \]
 because $|z|\le4+r_u$ on $\gamma_u$.  Therefore
\[
 \|F_u(D-z)^{-1}\|_{0\to0}\le a(g)K_1\le q(g).
\]
For $q(g)<1$, Yarotsky's Neumann series~(15) converges in
$\mathcal X_0^-$ and yields
\[
 \|(D+F_u-z)^{-1}-(D-z)^{-1}\|_{0\to0}
 \le K_0\frac{q(g)}{1-q(g)}.
\]
 The length of $\gamma_u$ is $2\pi r_u$, and
 $r_uK_0\le1$.  Contour integration therefore proves
 \eqref{eq:g18-cmp1-kappa}.
\end{proof}

\begin{lemma}[Finite-support shell bridge]
\label{lem:g18-cmp1-finite-support-shell-bridge}
In the physical charge-odd coefficient space,
\begin{equation}
 \operatorname{spec}(D|_{\mathcal X_0^-})\cap(-\infty,5)=\{4\},
 \qquad
 \operatorname{Ran}R^-
 =\overline{\operatorname{span}}^{\,\|\cdot\|_0}
   \{w_{x,\alpha}:x\in\mathbb Z^3,\ \alpha=1,2,3\}.
 \label{eq:g18-cmp1-free-shell-span}
\end{equation}
\end{lemma}

\begin{proof}
Let $v_J$ be a finite-support gauge-invariant, charge-odd eigen-coefficient
of rescaled electric energy below $5$.  Choose a periodic side length $L$
larger than the occupied-edge bounding box and embed the support injectively
in the $L^3$ torus, extending by the cell vacuum.  No occupied path winds
  around the torus, so the embedded coefficient belongs to the trivial
  electric one-form sector.  The retained-shell theorem then forces its energy
  to be $4$ and its energy-$4$ component to be a linear combination of odd
  elementary plaquette vectors.  The cell Laplacians have discrete Casimir
  spectrum, and energy below $5$ bounds the number of occupied cells.  Thus
  this diagonal direct sum has no additional approximate spectrum below $5$.
  Finite coefficient families are a core for
$D$, and spectral truncation in each finite tensor factor followed by
truncation in $J$ is dense in $\mathcal X_0^-$.  Applying these truncations
to the spectral projection of $D$ proves
\eqref{eq:g18-cmp1-free-shell-span}.
\end{proof}

\begin{lemma}[Coefficient totality]
\label{lem:g18-cmp1-coefficient-totality-complete}
Assume $\kappa(g)<1$.  Then
\begin{equation}
 \operatorname{Ran}\mathcal P_u^-
 =\overline{\mathcal P_u^-\operatorname{Ran}R^-}^{\,\|\cdot\|_0}.
 \label{eq:g18-cmp1-two-projection-totality}
\end{equation}
Using the finite-support retained-shell classification,
\begin{equation}
 \operatorname{Ran}\mathcal P_u^-
 =\overline{\operatorname{span}}^{\,\|\cdot\|_0}
 \{\mathcal P_u^-w_{x,\alpha}:
 x\in\mathbb Z^3,\ \alpha=1,2,3\}.
 \label{eq:g18-cmp1-seed-totality}
\end{equation}
\end{lemma}

\begin{proof}
Write $P=\mathcal P_u^-$, $R=R^-$, and $Q=1-R$.  Both $R$ and $Q$ are
blockwise orthogonal projections and hence contractions in the direct-sum
$\ell^1$ norm.  For $v\in\mathcal X_0^-$ set
\[
 y_0=Qv,
 \qquad y_{n+1}=QPy_n.
\]
Every $y_n$ lies in $\operatorname{Ran}Q$, and hence
\begin{equation}
 \|y_{n+1}\|_0
 =\|Q(P-R)y_n\|_0
 \le\kappa(g)\|y_n\|_0.
 \label{eq:g18-cmp1-geometric-irrelevant}
\end{equation}
Idempotence of $P$ gives the exact identity
\begin{equation}
 Py_n=P^2y_n=PRPy_n+PQPy_n=PRPy_n+Py_{n+1}.
 \label{eq:g18-cmp1-iteration-identity}
\end{equation}
Consequently, for every $N\ge0$,
\begin{equation}
 Pv=PRv+\sum_{n=0}^{N}PRPy_n+Py_{N+1}.
 \label{eq:g18-cmp1-iteration-finite}
\end{equation}
The last term tends to zero by
\eqref{eq:g18-cmp1-geometric-irrelevant}.  Every preceding term belongs to
$P\operatorname{Ran}R$, which proves~\eqref{eq:g18-cmp1-two-projection-totality}.
The reverse inclusion is immediate because $P$ is a projection.

 Lemma~\ref{lem:g18-cmp1-finite-support-shell-bridge} says that
 $\operatorname{Ran}R^-$ is the $\ell^1$-closed span of the
 $w_{x,\alpha}$.
Applying the bounded operator $P$ and using
\eqref{eq:g18-cmp1-two-projection-totality} proves
\eqref{eq:g18-cmp1-seed-totality}.  This is the complete version of the
iteration compressed into equations~(21)--(24) of
\cite{Yarotsky2004QP}; no one-site nondegeneracy assumption is used.
\end{proof}

\begin{lemma}[Coefficient-to-GNS transport]
\label{lem:g18-cmp1-gns-transport-complete}
 Let $\mathcal H_{u,\mathrm{phys}}^-$ be the closed local gauge-invariant,
 charge-odd cyclic subspace in the GNS representation of the thermodynamic
 ground state.  Put
 \begin{equation}
  \widehat H_\infty^-(u)=\frac32H_\infty^-(u),
  \label{eq:g18-cmp1-rescaled-H-infty}
 \end{equation}
 and let $P_{F,\infty}^-(u)$ be its shell Riesz projection on
 $\gamma_u$.  This is the same projection as the shell Riesz projection of
 $H_\infty^-(u)$, because multiplication of the Hamiltonian and contour by
 the same positive scalar does not change a spectral projection.
Then the coefficient map extends to a contraction
\begin{equation}
 \Gamma_u:\mathcal X_0^-\longrightarrow
 \mathcal H_{u,\mathrm{phys}}^-,
 \qquad
 \Gamma_uv=\pi_u\!\left(\sum_J\widehat v_J\right)\Omega_u,
 \label{eq:g18-cmp1-Gamma}
\end{equation}
and satisfies
\begin{equation}
 P_{F,\infty}^-(u)\Gamma_u
 =\Gamma_u\mathcal P_u^-.
 \label{eq:g18-cmp1-riesz-intertwining}
\end{equation}
Moreover,
\begin{equation}
 \operatorname{Ran}P_{F,\infty}^-(u)
 =\overline{\operatorname{span}}
 \{P_{F,\infty}^-(u)
       \pi_u(\widehat w_{x,\alpha})\Omega_u\}.
 \label{eq:g18-cmp1-gns-totality}
\end{equation}
No injectivity assumption on $\Gamma_u$ is required.
\end{lemma}

\begin{proof}
For a finite coefficient family,
\[
 \|\Gamma_uv\|
 \le\sum_J\|\pi_u(\widehat v_J)\Omega_u\|
 \le\sum_J\|\widehat v_J\|
 =\sum_J\|v_J\|=\|v\|_0.
\]
Thus the series in~\eqref{eq:g18-cmp1-Gamma} converges in operator norm and
$\Gamma_u$ extends by continuity to $\mathcal X_0^-$.  Exact-support
symmetry, Lemma~\ref{lem:g18-cmp1-exact-support-symmetry}, puts its range in
the physical odd sector.

 Yarotsky's equation~(18), in the present gap-one normalization, gives on the
 finite coefficient core
 \begin{equation}
  \widehat H_\infty^-(u)\Gamma_uv=\Gamma_u(D+F_u)v.
 \label{eq:g18-cmp1-generator-intertwining}
\end{equation}
Finite exact-support and local spectral truncations approximate every
$v\in\mathcal X_1^-$ in the $D$-graph norm.  The relative bound
\eqref{eq:g18-cmp1-relative-bound}, boundedness of $\Gamma_u$, and closedness
 of $\widehat H_\infty^-(u)$ therefore extend
\eqref{eq:g18-cmp1-generator-intertwining} to all of
 $\mathcal X_1^-$.  If $z$ lies on $\gamma_u$ and $a\in\mathcal X_0^-$,
apply this equality to
$v=(D+F_u-z)^{-1}a$ to obtain
\begin{equation}
  (\widehat H_\infty^-(u)-z)^{-1}\Gamma_ua
 =\Gamma_u(D+F_u-z)^{-1}a.
 \label{eq:g18-cmp1-resolvent-intertwining}
\end{equation}
Contour integration proves~\eqref{eq:g18-cmp1-riesz-intertwining}.

 It remains only to justify density in the correct sector.  The paragraph
 preceding equation~(18) of \cite{Yarotsky2004QP} states that finite linear
 combinations of creator vectors $\pi_u(\widehat v_J)\Omega_u$ are dense in
 the GNS Hilbert space.  To restrict this dense set to the physical odd
 sector, average a finite creator family over
the compact product of gauge groups at the finitely many vertices incident
to its support, and then apply $(1-C)/2$.  By
Lemma~\ref{lem:g18-cmp1-exact-support-symmetry}, this operation neither
enlarges exact support nor leaves the finite coefficient core.  Since the
corresponding Hilbert-space average is an orthogonal projection that fixes
every physical odd vector, the averaged creator families are dense in
$\mathcal H_{u,\mathrm{phys}}^-$.  Hence
\[
 \operatorname{Ran}P_{F,\infty}^-(u)
 =\overline{P_{F,\infty}^-(u)\Gamma_u\mathcal c_{00}^-}
 =\overline{\Gamma_u\mathcal P_u^-\mathcal c_{00}^-}.
\]
Now apply Lemma~\ref{lem:g18-cmp1-coefficient-totality-complete} and the
boundedness of $\Gamma_u$ to obtain
\eqref{eq:g18-cmp1-gns-totality}.
\end{proof}

\begin{corollary}[Uniform nonvanishing of each canonical seed]
\label{cor:g18-cmp1-seed-nonzero-complete}
For all sufficiently small $g$, every vector
\[
 \psi_{x,\alpha}
 =P_{F,\infty}^-(u)\pi_u(\widehat w_{x,\alpha})\Omega_u
\]
is nonzero, uniformly in $x$ and $\alpha$.
\end{corollary}

\begin{proof}
 The normalized creator has three-cell support $S_s$ and obeys
 $\widehat w_s^{\,*}\widehat w_s
 =p_{S_s}:=|\Omega_{S_s,0}\rangle\langle\Omega_{S_s,0}|\otimes I$.
 Here is the needed local product-vacuum estimate.  In a periodic volume,
 the product vacuum is a trial state, while the sum of interaction terms is
 bounded below by $-|\Lambda|g$.  Hence
 \[
  E_\Lambda(u)\le |\Lambda|g,
  \qquad
  E_\Lambda(u)\ge
  \sum_{x\in\Lambda}\langle\widehat h_x\rangle_{\Lambda,u}
       -|\Lambda|g .
 \]
 Translation invariance gives
 $\langle\widehat h_x\rangle_{\Lambda,u}\le2g$.  Since
 $1-p_x\le\widehat h_x$ and
 $1-p_{S_s}\le\sum_{x\in S_s}(1-p_x)$, every periodic subsequential limit
 obeys
 \[
  1-\omega_u(p_{S_s})\le2|S_s|g=6g.
 \]
 Yarotsky uniqueness identifies that limit with $\omega_u$.
Since $R^-w_s=w_s$, Lemma~\ref{lem:g18-cmp1-close-projection-complete} and
the contraction property of $\Gamma_u$ imply
\begin{align*}
 \|\psi_s\|
 &=\|\Gamma_u\mathcal P_u^-w_s\|\\
 &\ge \|\Gamma_uw_s\|
       -\|\Gamma_u(\mathcal P_u^--R^-)w_s\|\\
 &\ge \sqrt{1-6g}-\kappa(g).
\end{align*}
The last quantity is positive after decreasing the small-coupling interval.
\end{proof}

\begin{remark}[Scope]
\label{rem:g18-cmp1-scope-complete}
Lemmas~\ref{lem:g18-cmp1-close-projection-complete}--
\ref{lem:g18-cmp1-gns-transport-complete} close CMP(1) for the complete
three-component shell at fixed lattice spacing.  They do not identify an
individual internal dispersion sheet, an isotropic Wilson transfer-matrix
particle, or a continuum Yang--Mills particle.
\end{remark}
